\section*{Chapitre \ref{ch:s2:gamma-scalping-and-event-volatility} --- Gamma scalping et volatilité d'événement}

\begin{solution}{exo:s2:gamma-scalping-and-event-volatility:1}
$\sqrt{2/\pi}\times\sqrt{0.0025} = \sqrt{2/\pi}\times 0.05 = 3.99~\%$.
\end{solution}

\begin{solution}{exo:s2:gamma-scalping-and-event-volatility:2}
$\sigma^2 = 0.09 + 0.0025/(10/252)$, donc $\sigma = 39.1~\%$.
\end{solution}

\begin{solution}{exo:s2:gamma-scalping-and-event-volatility:3}
L'annonce intervient hors des heures de marché et le cours saute à l'ouverture~: aucune transaction n'a lieu entre la dernière couverture avant le rapport et le
premier cours qui suit, si bien que tout le mouvement est subi au delta détenu pendant la nuit.
\end{solution}

\begin{solution}{exo:s2:gamma-scalping-and-event-volatility:4}
La plupart des straddles versent la prime de variance les jours de diffusion, un risque que les vendeurs sont payés pour supporter. Autour des annonces, les vendeurs sous-estiment l'incertitude liée
à l'événement, davantage pour les petites sociétés volatiles et coûteuses à négocier~; la prime d'événement joue dans l'autre sens pendant quelques jours.
\end{solution}

\begin{solution}{exo:s2:gamma-scalping-and-event-volatility:5}
Il se couvre lorsque le delta a bougé, c'est-à-dire quand la trajectoire a livré une variance qui mérite d'être encaissée, et non
pendant les intervalles calmes, où une couverture s'exécute pour peu de bénéfice.
\end{solution}

\begin{solution}{exo:s2:gamma-scalping-and-event-volatility:6}
Le mouvement supplémentaire lié à l'événement a un écart-type d'environ 0.41~\%, contre environ 0.7~\% pour les mouvements quotidiens d'un indice~; un saut de cours sur résultats d'une action est souvent
plusieurs fois son mouvement quotidien. La variance d'événement intégrée dans les options vendues sur indice est faible, et la marge de l'opération est d'autant plus mince.
\end{solution}

\begin{solution}{exo:s2:gamma-scalping-and-event-volatility:7}
Juste en moyenne, le straddle couvert chaque jour perd 2.2~\% par événement~; à 15~\% trop bas, il perd 0.1~\%~; à 30~\% trop bas, il gagne 2.1~\%. Le seuil de rentabilité se situe un peu
au-delà de 15~\% trop bas. Le résultat n'est pas nul, car le straddle paie trois jours de theta, à une volatilité qui porte la prime de diffusion, et le coût aller-retour
de l'option ainsi que les couvertures avant l'événement ne rapportent rien.
\end{solution}

\begin{solution}{exo:s2:gamma-scalping-and-event-volatility:8}
Le 1.04 suppose 200 événements indépendants par an. Les saisons de résultats regroupent les publications sur quelques semaines, et les événements partagent des mouvements de marché et de secteur, si bien que le
nombre effectif de paris indépendants est plus faible, et le ratio de Sharpe plus bas~; de plus, l'écart-type d'un événement atteint 29~\% de la prime. Il faut dimensionner sur
un scénario de stress d'une mauvaise saison, et non sur le ratio de Sharpe.
\end{solution}

\begin{solution}{pb:s2:gamma-scalping-and-event-volatility:1}
\begin{enumerate}
\item $\tfrac12\Gamma S^2(\sigma_r^2 - \sigma_i^2)\,\mathrm{d}t$~: le gamma multiplié par la variance livrée, moins la variance payée par le theta.
\item L'espérance est fixée par la variance implicite comparée à la variance réalisée~; la couverture modifie seulement la manière dont la variance de la trajectoire est collectée, d'où le bruit.
\item $-12.1~\%$ sans couverture, à $-14.1~\%$ avec une couverture horaire, avec des écarts-types descendant de 62.4~\% à 6.7~\%, et des coûts de couverture atteignant 2.1~\%.
\item Une règle de bande couvre lorsque le delta a bougé de plus d'un seuil~; une règle d'horloge, à heures fixes. Les bandes réduisent le bruit pour un même coût.
\item Deux échéances après l'événement~: $\sigma^2(\tau)\tau = \sigma_d^2\tau + E$, deux équations en $\sigma_d$ et $E$.
\item Le \omterm{def:s2:gamma-scalping-and-event-volatility:move}{mouvement implicite} est le rendement absolu espéré de l'événement, $\sqrt{2/\pi}\sqrt{E}$~; le \omterm{def:s2:gamma-scalping-and-event-volatility:straddle}{straddle d'événement} s'achète avant l'événement et se vend après.
\item L'incertitude anticipée est forte, variable dans le temps et informative (Dubinsky et coauteurs)~; les \omterm{def:s2:gamma-scalping-and-event-volatility:straddle}{straddles d'événement} ont rapporté en moyenne 3.34~\% à partir de trois
jours avant jusqu'à la date d'annonce (Gao, Xing et Zhang).
\item La variance d'événement implicite vaut la vraie multipliée par 0.70 et par une erreur lognormale~; le \omterm{def:s2:gamma-scalping-and-event-volatility:move}{mouvement implicite} est en moyenne trop bas et erroné
d'un montant variable.
\item La baisse de la volatilité implicite une fois l'événement passé~; une médiane de 36.6~\% à 32.0~\%, soit 8.6~\%.
\item Rendement absolu moyen de 0.83~\% contre 0.71~\%~; le VIX9D a baissé lors de 65.6~\% des jours de communiqué, de 3.05~\% en moyenne.
\item La variance de l'événement représente une part plus faible de la variance à 30 jours que de celle à 9 jours.
\item Le vega perdu le jour de l'événement, quel que soit le comportement de l'action.
\item \textbf{Résultat nommé.} Le \omterm{def:s2:gamma-scalping-and-event-volatility:straddle}{straddle d'événement} rapporte 2.1~\% de sa prime par événement, avec un écart-type de 29~\%, de $+8.0\%$ dans le cinquième des
événements les plus sous-évalués à $-4.5~\%$ dans le cinquième évalué à peu près correctement~; la couverture quotidienne maximise le rendement après coûts (1.04 par an pour 200 événements).
\item 8.0, 2.7, 3.9, 0.5 et $-4.5~\%$, du cinquième le plus sous-évalué au cinquième le moins sous-évalué.
\item Le mouvement effectivement observé est un seul tirage de sa distribution.
\item Environ 15~\% de variance sous-estimée suffisent tout juste à l'équilibre~; 30~\% rapportent 2.1~\%.
\item Événements indépendants.
\item Cotations des options à l'entrée et à la sortie, dates d'annonce telles qu'elles étaient connues à l'époque, \omterm{def:s2:gamma-scalping-and-event-volatility:crush}{effondrement de la volatilité} à la sortie, coûts, saisons qui regroupent les événements.
\item La vente de volatilité après l'événement.
\item La probabilité que l'événement fasse bouger l'action davantage que ne l'anticipe le marché.
\end{enumerate}
\end{solution}

\begin{solution}{iq:s2:gamma-scalping-and-event-volatility:1}
Détenir des options et les couvrir en delta, de sorte que les couvertures encaissent la variance de la trajectoire pendant que le theta paie la variance implicite~; la stratégie gagne de l'argent lorsque
la volatilité réalisée dépasse la volatilité implicite d'un montant supérieur aux coûts de couverture.
\end{solution}

\begin{solution}{iq:s2:gamma-scalping-and-event-volatility:2}
À partir du straddle à la monnaie dont l'échéance tombe juste après le rapport, exprimé en proportion du cours de l'action~; ou, mieux, extraire la variance d'événement
de deux échéances situées après le rapport et prendre $\sqrt{2/\pi}\sqrt{E}$.
\end{solution}

\begin{solution}{iq:s2:gamma-scalping-and-event-volatility:3}
Comparer les \omterm{def:s2:gamma-scalping-and-event-volatility:move}{mouvements implicites} aux mouvements réalisés le jour de l'événement, sur de nombreuses publications de résultats, en contrôlant la taille et la volatilité~;
négocier les \omterm{def:s2:gamma-scalping-and-event-volatility:straddle}{straddles d'événement} aux cotations affichées, coûts compris~; vérifier la stabilité d'une année sur l'autre et selon le type de société.
\end{solution}

\begin{solution}{iq:s2:gamma-scalping-and-event-volatility:4}
Des sauts de cours plus grands que ceux implicites sur de nombreuses valeurs à la fois (une saison de surprises), des mouvements sectoriels et des sauts corrélés~; tester en stress
les plus grandes saisons passées et un saut conjoint sur un secteur.
\end{solution}

\begin{solution}{iq:s2:gamma-scalping-and-event-volatility:5}
Positions en temps réel, cours, coûts et bande par portefeuille~; à l'ouverture, les cours sautent et les premières cotations sont larges, si bien qu'une couverture sur la première transaction peut
s'exécuter à un mauvais prix~; le moteur a besoin de règles pour les enchères et les cotations périmées.
\end{solution}

\begin{solution}{iq:s2:gamma-scalping-and-event-volatility:6}
$\sigma_d^2 = (\sigma_2^2\tau_2 - \sigma_1^2\tau_1)/(\tau_2 - \tau_1)$ et $E = \sigma_1^2\tau_1 - \sigma_d^2\tau_1$. Une échéance antérieure à l'événement
ne porte aucune variance d'événement, de sorte que la formule traiterait sa volatilité plus faible comme diffusive et que le $E$ extrait serait faux.
\end{solution}
